\specialchapter{BIBLIOGRAPHY}

{[}1{]} R. ARENS - "On a theorem of Gelfand and Neumark", Proc. Natl. Acad. Sci. USA 32 (1946), p. 237-239, Zbl 0060.27101.

{[}2{]} M. F. ATIYAH and I. M. SINGER -- "The index of elliptic operators on compact manifolds", Bull. Amer. Math. Soc. 69 (1963), p. 422--433, Zbl 0118.31203; Michael Atiyah Collected Works, vol. 3, Clarendon Press, 1988, p. 13--24.

{[}3{]} F. V. ATKINSON -- "Normal solvability of linear equations in normed spaces", Mat. Sb., Nov. Ser. 28 (1951), p. 3--14, Zbl 0042.12001 (in Russian).

{[}4{]} S. BANACH - Theory of linear operations, Warsaw, 1932, Zbl 0005.20901.

{[}5{]} V. BARGMANN - "Irreducible unitary representations of the Lorentz group", Ann. of Math. (2) 48 (1947), p. 568-640, Zbl 0045.38801.

{[}6{]} S. BOCHNER -- "Contributions to the theory of almost periodic functions. II: Functions of several variables", Math. Ann. 96 (1926), p. 383--409, JfM 52.0265.01; Collected papers, vol. 2, Amer. Math. Soc., 1992, p. 79--106.

{[}7{]} \_\_\_\_ Lectures on Fourier integrals, Leipzig, 1932, Zbl 0006.11001.

{[}8{]} \_\_\_\_ "Spectral decomposition of linear families of unitary operators", Sitzungsber. Preuß. Akad. Wiss., Phys.-Math. Kl. (1933), p. 371--376, Zbl 0007.16603; Collected papers, vol. 1, Amer. Math. Soc., 1992, p. 579--586.

{[}9{]} H. BOHR - "On the theory of almost periodic functions. I. A generalization of the theory of Fourier series", Acta Math. 45 (1924), p. 29-127, JfM 50.0196.01.

{[}10{]} J. W. CALKIN - "Two-sided ideals and congruences in the ring of bounded operators in Hilbert space", Ann. of Math. (2) 42 (1941), p. 839-873, Zbl 0063.00692.

{[}11{]} T. CARLEMAN - On singular integral equations with real and symmetric kernel, Uppsala, 1923, JfM 49.0272.01.

{[}12{]} H. CARTAN and R. GODEMENT -- "Duality theory and harmonic analysis in locally compact abelian groups", Ann. Sci. Éc. Norm. Supér. (3) 64 (1947), p. 79--99, Zbl 0033.18801; Henri Cartan's Works, vol. III, Springer, 1979, p. 1203--1223.

{[}13{]} H. CARTAN and J-P. SERRE -- "A finiteness theorem concerning compact analytic manifolds", C. R. Acad. Sci., Paris 237 (1953), p. 128--130, Zbl 0050.17701; Henri Cartan's Works, vol. II, Springer, 1979, p. 684--686.

{[}14{]} C. CHEVALLEY - "Generalization of class field theory for infinite extensions", J. Math. Pures Appl. (9) 15 (1936), p. 359-371, Zbl 0015.15101.

{[}15{]} J. DIEUDONNÉ -- "On homomorphisms of normed spaces", Bull. Sci. Math. (2) 67 (1943), p. 72--84, Zbl 0028.23301; Selected mathematical works, Hermann, 1981, p. 268--281.

{[}16{]} J. DIEUDONNÉ - History of functional analysis, North-Holland Publ. Co., 1981, Zbl 0478.46001.

{[}17{]} P. L. DIRICHLET -- "Proof of the theorem that every unbounded arithmetic progression, whose first term and difference are integers without common factor, contains infinitely many primes", Abh. Preuß. Akad. Wiss. 8 (1837), p. 45--81; Works, vol. 1, Reimer (Berlin), 1889, p. 315--356.

{[}18{]} P. ENFLO - "A counterexample to the approximation problem in Banach spaces", Acta Math. 130 (1973), p. 309-317, Zbl 0267.46012.

{[}19{]} \_\_\_\_ "On the invariant subspace problem for Banach spaces", Acta Math. 158 (1987), no. 3-4, p. 213--313, Zbl 0663.47003.

{[}20{]} J. M. G. FELL - "The dual spaces of C*-algebras", Trans. Amer. Math. Soc. 94 (1960), p. 365-403, Zbl 0090.32803.

{[}21{]} M. FUKAMIYA - "On a theorem of Gelfand and Neumark and the B*-algebra", Kumamoto J. Sci., Math. 1 (1952), p. 17-22, Zbl 0049.08605.

{[}22{]} C. F. GAUSS - Disquisitiones Arithmeticae, Göttingen, 1801.

{[}23{]} I. M. GELFAND - Collected papers, 3 vol., Springer, 1987.

{[}24{]} I. M. GELFAND and M. A. NAIMARK -- "Normed rings with involution and their representations", Izv. Akad. Nauk SSSR, Ser. Mat. 12 (1948), p. 445--480, Zbl 0031.03403; I. M. Gelfand Collected papers, vol. I, Springer, 1987, p. 364--400 (in Russian).

{[}25{]} R. GODEMENT -- "On the orthogonality relations of V. Bargmann. II. General proof", C. R. Acad. Sci. Paris 225 (1947), p. 657--659, Zbl 0029.19907.

{[}26{]} I. C. GOHBERG -- "On linear equations in normed spaces", Dokl. Akad. Nauk SSSR, n. Ser. 76 (1951), p. 477--480, Zbl 0042.11903 (in Russian).

{[}27{]} P. R. HALMOS - "What does the spectral theorem say?", Amer. Math. Monthly 70 (1963), p. 241-247, Zbl 0132.35606; Selecta, vol. II, Springer, 1983, p. 59-66.

{[}28{]} W. HEISENBERG -- "On a modification of the formal rules of quantum theory in the problem of the anomalous Zeeman effects", Z. Phys. 26 (1924), p. 291--307, JfM 51.0740.07; Collected Works: Scientific original papers, vol. 1, Springer, 1985, p. 306--325.

{[}29{]} E. HELLINGER and O. TOEPLITZ -- "Integral equations and equations with infinitely many unknowns", Enzyklopädie der mathematischen Wissenschaften, art. II C 13, Leipzig (Teubner), 1927, JfM 53.0350.01.

{[}30{]} E. HILB - "On integral representation of arbitrary functions", Math. Ann. 66 (1909), p. 1-66, JfM 39.0380.03.

{[}31{]} D. HILBERT - "Foundations of a general theory of linear integral equations. First communication.", Nachr. Ges. Wiss. Göttingen, Math.-Phys. Kl. (1904), p. 49-91, JfM 35.0378.02.

{[}32{]} \_\_\_\_ "Foundations of a general theory of linear integral equations. Fourth communication.", Nachr. Ges. Wiss. Göttingen, Math.-Phys. Kl. (1906), p. 157--227, JfM 37.0351.03.

{[}33{]} T. H. HILDEBRANDT - "On completely continuous linear transformations", Acta Math. 51 (1928), p. 311-318, JfM 54.0427.03.

{[}34{]} E. HILLE - Functional analysis and semi-groups, Colloq. Publ., vol. 31, Amer. Math. Soc., 1948, Zbl 0033.06501.

{[}35{]} M. JAMMER - The conceptual development of quantum mechanics, McGraw-Hill, 1966.

{[}36{]} M. KAC - "Can one hear the shape of a drum?", Am. Math. Mon. 73 (1966), p. 1-23, Zbl 0139.05603.

{[}37{]} E. R. VAN KAMPEN - "Locally bicompact abelian groups and their character groups", Ann. of Math. (2) 36 (1935), p. 448-463, Zbl 0011.39401.

{[}38{]} \_\_\_\_ "Almost periodic functions and compact groups", Ann. of Math. (2) 37 (1936), p. 78--91, Zbl 0014.01702.

{[}39{]} P. LÉVY - "On the absolute convergence of Fourier series", Compos. Math. 1 (1934), p. 1-14, Zbl 0008.31201.

{[}40{]} V. I. LOMONOSOV - "Invariant subspaces for the family of operators which commute with a completely continuous operator", Funct. Anal. Appl. 7 (1974), p. 213-214, Zbl 0293.47003.

{[}41{]} L. H. LOOMIS - An introduction to abstract harmonic analysis, Van Nostrand Co., 1953, Zbl 0052.11701.

{[}42{]} G. W. MACKEY - "Imprimitivity for representations of locally compact groups. I", Proc. Nat. Acad. Sci. U.S.A. 35 (1949), p. 537-545, Zbl 0035.06901.

{[}43{]} M. MATHIAS - "On positive Fourier integrals", Math. Z. 16 (1923), p. 103--125, JfM 49.0207.03.

{[}44{]} S. MAZUR - "On linear rings", C. R. Acad. Sci., Paris 207 (1938), p. 1025-1027, Zbl 0020.20101.

{[}45{]} J. MERCER - "Functions of positive and negative type, and their connection with the theory of integral equations", Philos. Trans. R. Soc. Lond., Ser. A 209 (1909), p. 415-446, JfM 40.0408.02.

{[}46{]} S. G. MIKHLIN - "On the resolution of linear equations in Hilbert spaces", Dokl. Akad. Nauk SSSR (N. S.) 57 (1947), p. 11-12, Zbl 0029.05101 (in Russian).

{[}47{]} F. J. MURRAY and J. VON NEUMANN - "On rings of operators", Ann. of Math. (2) 37 (1936), p. 116-229, Zbl 0014.16101; John von Neumann Collected works, vol. III, Pergamon Press, 1961, p. 6-120.

{[}48{]} M. NAGUMO - "Some analytic investigations in linear metric rings", Jpn. J. Math. 13 (1936), p. 61-80, Zbl 3024132.

{[}49{]} J. VON NEUMANN -- "Mathematical foundations of quantum mechanics", Göttingen Nachrichten 1927 (1927), p. 1--57, JfM 53.0848.03; Collected works, vol. I, Pergamon Press, 1961, p. 151--207.

{[}50{]} \_\_\_\_ "General eigenvalue theory of Hermitian functional operators", Math. Ann. 102 (1929), p. 49--131, JfM 55.0824.02; Collected works, vol. II, Pergamon Press, 1961, p. 3--85.

{[}51{]} \_\_\_\_ "On the algebra of functional operations and the theory of normal operators", Math. Ann. 102 (1929), p. 370--427, JfM 55.0825.02; Collected works, vol. II, Pergamon Press, 1961, p. 86--143.

{[}52{]} \_\_\_\_ "Almost periodic functions in a group. I", Trans. Amer. Math. Soc. 36 (1934), p. 445--492, Zbl 0009.34902; Collected works, vol. II, Pergamon Press, 1961, p. 454--501.

{[}53{]} F. NOETHER -- "On a class of singular integral equations", Math. Ann. 82 (1920), p. 42--63, JfM 47.0369.02.

{[}54{]} F. PETER and H. WEYL -- "The completeness of the primitive representations of a closed continuous group", Math. Ann. 97 (1927), p. 737--755, JfM 53.0387.02; Hermann Weyl Gesammelte Abhandlungen, Bd. III, Springer, 1968, p. 58--75.

{[}55{]} L. PONTRYAGIN -- "The general duality theorem for closed sets", Verhandlungen Kongreß Zürich, 2, p. 195-197, 1932, JfM 58.0646.08.

{[}56{]} \_\_\_\_ "Almost periodic functions and analysis situs", C. R. Acad. Sci., Paris 196 (1933), p. 1201--1203, Zbl 0006.42801.

{[}57{]} \_\_\_\_ "The general topological theorem of duality for closed sets", Ann. of Math. (2) 35 (1934), no. 4, p. 904--914, Zbl 0010.18004.

{[}58{]} \_\_\_\_ "The theory of topological commutative groups", Ann. of Math. (2) 35 (1934), p. 361--388, Zbl 0009.15601.

{[}59{]} F. RIESZ - Les systèmes d'équations linéaires à une infinité d'inconnues, Gauthier-Villars, 1913, JfM 44.0401.01.

{[}60{]} \_\_\_\_ "On linear functional equations", Acta Math. 41 (1916), p. 71--98, JfM 46.0635.01 ; Œuvres complètes, t. II, Gauthier-Villars, 1960, p. 1053--1080.

{[}61{]} \_\_\_\_ "On theorems of Stone and Bochner", Acta Litt. Sci. Szeged 6 (1933), p. 184--198, Zbl 0007.30903; Œuvres complètes, t. II, Gauthier-Villars, 1960, p. 1135--1149.

{[}62{]} J. SCHATZ - "Review of the article On a theorem of Gelfand and Neumark and the B*-algebra by M. Fukamiya", Math. Reviews 14 (1952), p. 884.

{[}63{]} J. SCHAUDER - "On linear, completely continuous functional operations", Studia Math. 2 (1930), p. 183-196, JfM 56.0354.01.

{[}64{]} A. SCHUSTER -- "The genesis of spectra", Report on the fifty-second meeting of the British association for the advancement of Science, J. Murray (London), 1882, p. 120--143.

{[}65{]} L. SCHWARTZ - "Generalization of the notion of function, of derivation, of Fourier transformation and mathematical and physical applications", Ann. Univ. Grenoble, Sect. Sci. Math. Phys., n. Ser. 21 (1946), p. 57-74, Zbl 0060.27504; Œuvres scientifiques, t. I, Soc. Math. France, 2011, p. 61-88.

{[}66{]} \_\_\_\_ "Homomorphisms and completely continuous mappings", C. R. Acad. Sci., Paris 236 (1953), p. 2472--2473, Zbl 0050.33301 ; Œuvres scientifiques, t. I, Soc. Math. France, 2011, p. 345--346.

{[}67{]} I. E. SEGAL - "The group algebra of a locally compact group", Trans. Amer. Math. Soc. 61 (1947), p. 69-105, Zbl 0032.02901.

{[}68{]} G. E. SHILOV -- "On the decomposition of a commutative normed ring into a direct sum of ideals", Mat. Sb., Nov. Ser. 32 (1953), p. 353--364, Zbl 0052.34203 (in Russian).

{[}69{]} S. SOBOLEV - "New method to solve the Cauchy problem for normal hyperbolic linear equations", Rec. Math. Moscou, n. Ser. 1 (1936), p. 39-71, Zbl 0014.05902.

{[}70{]} S. W. P. STEEN - "An introduction to the theory of operators : I. Real operators. Modulus", Proc. Lond. Math. Soc. (2) 41 (1936), p. 361-392, Zbl 0014.16201.

{[}71{]} T.-J. STIELTJES - "Researches on continued fractions", Ann. Fac. Sci. Toulouse Sci. Math. Sci. Phys. 8 (1894), no. 4, p. J1-J122, JfM 25.0326.01.

{[}72{]} M. H. STONE - Linear transformations in Hilbert space and their applications to analysis, Colloq. Publ., Am. Math. Soc., vol. 15, Amer. Math. Soc., 1932, Zbl 0005.40003.

{[}73{]} \_\_\_\_ "On one-parameter unitary groups in Hilbert space", Ann. of Math. (2) 33 (1932), p. 643--648, Zbl 0005.16403.

{[}74{]} \_\_\_\_ "Applications of the theory of Boolean rings to general topology", Trans. Am. Math. Soc. 41 (1937), p. 375--481, Zbl 0017.13502.

{[}75{]} J. T. TATE — “Fourier analysis in number fields, and Hecke's zetafunctions”, Thesis, Princeton, 1950; Collected works, part I, Amer. Math. Soc., 2016, p. 1–43, Zbl 1407.01030.

{[}76{]} O. TOEPLITZ — “Das algebraische Analogon zu einem Satze von Fejér”, Math. Z. 2 (1918), p. 187–197, JfM 46.0157.02.

{[}77{]} H. WEBER — “Beweis des Satzes, dass jede eigentlich primitive quadratische Form unendlich viele Primzahlen darzustellen fähig ist.”, Math. Ann. 20 (1882), p. 301–330, JfM 14.0141.01.

{[}78{]} \_\_\_\_ “Theorie der Abel'schen Zahlkörper”, Acta Math. 9 (1887), p. 105–130, JfM 18.0055.04.

{[}79{]} A. WEIL — “Sur les fonctions presque périodiques”, C. R. Acad. Sci., Paris 200 (1935), p. 38–40, Zbl 0010.35303; Œuvres scientifiques, vol. I, Springer, 1979, p. 101–103.

{[}80{]} \_\_\_\_ L'intégration dans les groupes topologiques et ses applications, Actualités scientifiques et industrielles, vol. 869, Hermann, 1940, Zbl 0063.08195.

{[}81{]} H. WEYL — “Über beschränkte quadratische Formen, deren Differenz vollstetig ist”, Rend. Circ. Mat. Palermo 27 (1909), p. 373–392, JfM 40.0395.01; Gesammelte Abhandlungen, Bd. I, Springer, 1968, p. 175–194.

{[}82{]} \_\_\_\_ “Über gewöhnliche Differentialgleichungen mit Singularitäten und die zugehörigen Entwicklungen willkürlicher Funktionen”, Math. Ann. 68 (1910), p. 220–269, JfM 41.0343.01; Gesammelte Abhandlungen, Bd. I, Springer, 1968, p. 248–297.

{[}83{]} \_\_\_\_ “Theorie der Darstellung kontinuierlicher halbeinfacher Gruppen durch lineare Transformationen. I, II, III, Nachtrag”, Math. Z. 23 (1925), p. 271–309, JfM 51.0319.01; Gesammelte Abhandlungen, Bd. II, Springer, 1968, p. 543–647.

{[}84{]} \_\_\_\_ “Integralgleichungen und fastperiodische Funktionen”, Math. Ann. 97 (1926), p. 338–356, JfM 52.0260.02; Gesammelte Abhandlungen, Bd. III, Springer, 1968, p. 38–57.

{[}85{]} \_\_\_\_ “Quantenmechanik und Gruppentheorie”, Z. Phys. 46 (1927), p. 1–46, JfM 53.0848.02; Gesammelte Abhandlungen, Bd. III, Springer, 1968, p. 90–135.

{[}86{]} \_\_\_\_ Gruppentheorie und Quantenmechanik, S. Hirzel (Leipzig), 1928, Zbl 0537.22023.

{[}87{]} E. WEYR — “Zur Theorie der bilinearen Formen”, Monatsh. Math. Phys. 1 (1890), p. 163–236, JfM 22.0141.02.

{[}88{]} N. WIENER — “Tauberian theorems”, Ann. of Math. (2) 33 (1932), p. 1–100, Zbl 0004.05905; Selected Papers, MIT Press, 1964, p. 261–360.

{[}89{]} N. WIENER and R. E. A. C. PALEY — “Analytic properties of the characters of infinite abelian groups”, Verhandlungen Kongreß Zürich, 2, p. 95, 1932, JfM 58.0445.02.

{[}90{]} E. WIGNER — “On unitary representations of the inhomogeneous Lorentz group”, Ann. of Math. (2) 40 (1939), p. 149–204, Zbl 0020.29601; Collected works, part A, Springer, 1993, p. 334–389.

{[}91{]} W. WIRTINGER — “Beiträge zu Riemann's Integrationsmethode für hyperbolische Differentialgleichungen, und deren Anwendungen auf Schwingungsprobleme”, Math. Ann. 48 (1897), p. 365–389, JfM 27.0281.01.

{[}92{]} S. L. WORONOWICZ — “Unbounded elements affiliated with C*-algebras and noncompact quantum groups”, Comm. Math. Phys. 136 (1991), no. 2, p. 399–432, Zbl 0743.46080.

{[}93{]} B. YOOD — “Properties of linear transformations preserved under addition of a completely continuous transformation”, Duke Math. J. 18 (1951), p. 599–612, Zbl 0043.11901.

